% ==========================================================================
%  Chapter 71 -- Persistent Nonlinear Sub-Model Evolver and L-Shaped
%                Building Multi-Scale Seismic Analysis
% ==========================================================================

\chapter{Persistent Nonlinear Sub-Model Evolution and L-Shaped Building
         Multi-Scale Seismic Analysis}
\label{ch:nl-sub-model-evolver}

Previous chapters developed the kinematic reconstruction pipeline
(Chapter~\ref{ch:reconstruction}), the prismatic domain builder for
Hex8/20/27 meshes (Chapter~\ref{ch:ho-prismatic-domain}), and the
single-shot sub-model solver (\texttt{SubModelSolver}).  While the
single-shot solver is suitable for instantaneous snapshots of the
sub-model response, it \emph{discards the material state} after every
call---each solve starts from virgin concrete.  For an earthquake
time-history analysis where the KoBathe concrete model accumulates
cracks, plastic strain, and damage progressively, a persistent solver
is essential.

This chapter presents the \texttt{NonlinearSubModelEvolver}, a class
that maintains a continuum sub-model alive throughout the earthquake
and drives it step-by-step using direct PETSc SNES management.  It is
demonstrated through a 4-story L-shaped reinforced-concrete building
subjected to three-component earthquake excitation.


% =========================================================================
\section{Architecture}
\label{sec:nlevo-architecture}

The evolver differs from \texttt{SubModelSolver} in three critical
aspects:

\begin{enumerate}
  \item \textbf{Persistent \texttt{Model<>} instance.}  The model,
        its DMPlex mesh, and all material-point states are created once
        and reused across every global time step.  Crack records,
        plastic strains, and damage indices survive between steps.

  \item \textbf{Direct SNES management.}  Instead of using
        \texttt{NonlinearAnalysis::solve\_incremental()} (which
        resets the displacement vector $\mathbf{U}$ to zero before
        every solve), the evolver maintains its own \texttt{SNES},
        \texttt{Vec U}, \texttt{Vec R}, and \texttt{Mat J} objects.
        The converged $\mathbf{U}$ from step~$N$ serves as the
        initial guess for step~$N{+}1$.

  \item \textbf{Post-setup BC updates.}  After the model has been
        set up (constraints burned into the PetscSection), new
        boundary displacement values are written directly into the
        \texttt{imposed\_solution()} vector using PetscSection
        offsets, rather than calling \texttt{constrain\_node()}.
\end{enumerate}


% =========================================================================
\section{Solver Workflow}
\label{sec:nlevo-workflow}

\subsection{First Call (Model Creation + Incremental Loading)}

On the first invocation of \texttt{solve\_step()}:
\begin{enumerate}
  \item A \texttt{Model<ThreeDimensionalMaterial, SmallStrain, 3>} is
        created over the prismatic domain with KoBathe concrete.
  \item \texttt{constrain\_node()} is called for all face-A and face-B
        boundary nodes with the first set of displacement targets.
  \item \texttt{model->setup()} creates the DMPlex, the local/global
        vectors, and the PetscSection with constraint information.
  \item A persistent SNES is created with custom \texttt{FormResidual}
        and \texttt{FormJacobian} callbacks (identical to those in
        \texttt{NonlinearAnalysis} but owned by the evolver).
  \item The global displacement vector $\mathbf{U}$ is initialised
        to zero, and the target imposed solution is saved.
  \item An incremental loading loop scales the imposed displacements
        from $p{=}0$ to $p{=}1$ in $N$ sub-steps:
        \[
          \mathbf{u}_{\mathrm{imp}}^{(k)}
          = \frac{k}{N}\,\mathbf{u}_{\mathrm{target}},
          \quad k = 1, \ldots, N
        \]
        At each sub-step, \texttt{SNESSolve(snes, NULL, U)} is called
        from the current $\mathbf{U}$.  On convergence, the material
        state is committed; on divergence, it is reverted.
\end{enumerate}


\subsection{Subsequent Calls (BC Update + Single Newton Solve)}

For each subsequent global time step:
\begin{enumerate}
  \item New section kinematics $(\kappa_A, \kappa_B)$ are extracted
        from the beam element.
  \item Boundary displacements are recomputed via
        \texttt{compute\_boundary\_displacements()}.
  \item The \texttt{imposed\_solution()} vector is updated by
        obtaining PetscSection offsets for each constrained node
        (via \texttt{domain.node(nid).sieve\_id}) and writing the
        new displacement values directly.
  \item \texttt{SNESSolve(snes, NULL, U)} is called from the
        previously converged $\mathbf{U}$---no incremental loading
        is needed because the BC change between consecutive time
        steps is small.
  \item On convergence, material state is committed.
\end{enumerate}

The key insight is that the residual function computes the total
displacement as:
\[
  \mathbf{u}_{\mathrm{total}}
  = \underbrace{\mathbf{u}_{\mathrm{free}}}_{\text{from }\mathbf{U}}
  + \underbrace{\mathbf{u}_{\mathrm{imp}}}_{\text{imposed\_solution()}}
\]
When the imposed values change, the residual at the old
$\mathbf{U}$ reflects the new boundary conditions, and Newton
iterations adjust $\mathbf{U}$ to restore equilibrium.


% =========================================================================
\section{Crack Plane Visualisation}
\label{sec:nlevo-cracks}

After each converged step, the evolver collects
\texttt{InternalFieldSnapshot} data from every material point.
For each Gauss point with active cracks ($n_{\mathrm{crack}} \geq 1$),
a \texttt{CrackRecord} is stored containing:

\begin{itemize}
  \item Global coordinates of the Gauss point.
  \item Crack normal vectors $\mathbf{n}_1, \mathbf{n}_2, \mathbf{n}_3$.
  \item Crack opening strains $\varepsilon_1, \varepsilon_2, \varepsilon_3$.
  \item Open/closed status flags.
  \item Damage index $d$.
\end{itemize}

Each active crack is exported as a small \texttt{VTK\_QUAD} cell in an
unstructured grid (\texttt{.vtu}), oriented perpendicular to the crack
normal and centred at the Gauss point.  Cell data arrays store the
opening strain, normal vector, and crack state (open=1, closed=0).
These can be visualised in ParaView using the Glyph or Threshold
filters to highlight only open cracks or those exceeding a threshold
opening.

Three PVD time-series files are maintained per sub-model:
\begin{itemize}
  \item \texttt{nlsub\_<id>\_mesh.pvd} --- deformed mesh with
        displacement and material fields.
  \item \texttt{nlsub\_<id>\_gauss.pvd} --- Gauss-point cloud with
        stress, strain, and damage.
  \item \texttt{nlsub\_<id>\_cracks.pvd} --- crack plane quads.
\end{itemize}


% =========================================================================
\section{L-Shaped Building Example}
\label{sec:nlevo-lshaped}

\subsection{Structural Configuration}

A 4-storey reinforced-concrete moment frame with an L-shaped floor
plan is modelled.  The plan axes are:

\begin{center}
\begin{tabular}{ll}
  \toprule
  Direction & Grid coordinates \\
  \midrule
  $x$       & $0,\; 5,\; 10,\; 15$ m \\
  $y$       & $0,\; 4,\; 8$ m \\
  \bottomrule
\end{tabular}
\end{center}

The L-shape is achieved by cutting out grid cells
$(x_{\mathrm{start}}{=}2,\;x_{\mathrm{end}}{=}3,\;
  y_{\mathrm{start}}{=}1,\;y_{\mathrm{end}}{=}2)$
above all storeys, removing the upper-right quadrant of the plan.
Each storey has a height of 3.2~m; the total building height is 12.8~m.

Columns:\ $0.40 \times 0.40$~m, $f'_c = 28$~MPa.
Beams:\ $0.30 \times 0.50$~m, $f'_c = 25$~MPa.
Steel:\ $f_y = 420$~MPa (Menegotto--Pinto model).

Fiber sections are generated with a $6 \times 6$ grid for columns and
$5 \times 8$ for beams, with 4 reinforcing bars per section.


\subsection{Ground Motion}

Three components (NS, EW, UD) of the MYG004 Tsukidate record from the
2011 T\=ohoku earthquake are applied simultaneously:

\begin{center}
\begin{tabular}{lcl}
  \toprule
  Component & Direction & File \\
  \midrule
  NS & $x$ & \texttt{MYG0041103111446.NS} \\
  EW & $y$ & \texttt{MYG0041103111446.EW} \\
  UD & $z$ & \texttt{MYG0041103111446.UD} \\
  \bottomrule
\end{tabular}
\end{center}

Records are trimmed to 40--120~s and scaled by a factor of 0.5 to
maintain a moderate nonlinear response suitable for demonstrating the
multi-scale transition.


\subsection{Analysis Phases}

The analysis proceeds in two phases:

\paragraph{Phase 1: Global dynamic analysis with transition monitoring.}
A Newmark--$\beta$ (generalized-$\alpha$) time integration is performed
on the beam--column frame model.  Three observers are attached via a
\texttt{CompositeObserver}:
\begin{itemize}
  \item \texttt{DamageTracker} --- monitors curvature-based damage
        indices at every element.
  \item \texttt{FiberHysteresisRecorder} --- records $(\varepsilon,
        \sigma)$ histories for the 5 most-stressed fibres (segregated
        by material class: concrete and steel).
  \item \texttt{NodeRecorder} --- records displacement time histories
        at selected roof nodes.
\end{itemize}

A \texttt{TransitionDirector} halts Phase~1 when any element exceeds
a damage threshold (default $\varepsilon_y / E_s$).  At that instant,
the critical elements are identified and sub-models are created.

\paragraph{Phase 2: Multi-scale evolution.}
For each critical column element, a
\texttt{NonlinearSubModelEvolver} is instantiated with KoBathe
concrete ($f'_c = 28$~MPa).  The global dynamic analysis continues
step-by-step.  After each global step:
\begin{enumerate}
  \item Section kinematics are extracted from the beam element.
  \item Boundary displacements are computed for the sub-model faces.
  \item \texttt{solve\_step(t)} drives the sub-model from the previous
        converged state to the new BC target.
  \item Crack data is collected and VTK output is generated.
\end{enumerate}

This loop continues until the end of the earthquake record.  At
finalisation, PVD files are written for the global frame and for each
sub-model, and CSV files are exported for the recorders.


\subsection{Post-Processing and Visualisation}

The following output artefacts are generated:

\begin{description}
  \item[\texttt{roof\_displacement.csv}]
    Displacement time histories at selected roof nodes (columns:
    \texttt{time}, \texttt{node<id>\_dof<d>}).

  \item[\texttt{fiber\_hysteresis\_concrete.csv} /
        \texttt{fiber\_hysteresis\_steel.csv}]
    Stress--strain hysteresis loops for the 5 most-stressed fibres
    of each material class (columns:
    \texttt{time}, \texttt{e<e>\_gp<g>\_f<f>\_strain},
    \texttt{e<e>\_gp<g>\_f<f>\_stress}).

  \item[\texttt{frame\_NNNNNN.vtm}]
    VTK multi-block structural model at each output interval.

  \item[\texttt{nlsub\_<id>\_*.pvd}]
    PVD time series for each sub-model (mesh, Gauss cloud, crack planes).
\end{description}

A Python plotting script (\texttt{scripts/plot\_lshaped\_multiscale.py})
reads the CSV outputs and generates publication-quality figures:
\begin{itemize}
  \item Roof displacement time histories ($u_x$, $u_y$ vs.\ time).
  \item Roof orbit plot ($u_x$ vs.\ $u_y$).
  \item Concrete fibre hysteresis loops.
  \item Steel fibre hysteresis loops.
  \item Peak lateral displacement envelope.
\end{itemize}


% =========================================================================
\section{Implementation Notes}
\label{sec:nlevo-implementation}

\subsection{SNES Callback Registration}

The \texttt{FormResidual} and \texttt{FormJacobian} static methods
replicate the assembly logic of \texttt{NonlinearAnalysis} but are
defined as private static methods of \texttt{NonlinearSubModelEvolver}.
They receive a \texttt{Context\{model, f\_ext\}} pointer and perform:
\begin{enumerate}
  \item \texttt{DMGlobalToLocal(U)} to scatter free-DOF values.
  \item Add \texttt{imposed\_solution()} to obtain total displacement.
  \item Iterate model elements, calling
        \texttt{compute\_internal\_forces()} (residual) or
        \texttt{inject\_tangent\_stiffness()} (Jacobian).
  \item \texttt{DMLocalToGlobal} to assemble the residual.
\end{enumerate}

This avoids the \texttt{VecSet(U, 0.0)} that
\texttt{NonlinearAnalysis::solve\_incremental()} performs, enabling
Newton iterations to start from the previously converged state.


\subsection{Boundary Condition Updates After Setup}

After model setup, the constraint structure is encoded in the DMPlex
PetscSection and cannot be modified.  However, the \emph{values} of
constrained DOFs live in the \texttt{imposed\_solution()} vector
(a local sequential Vec).  The evolver writes new values at the
PetscSection offsets for each boundary node:

\begin{lstlisting}[language=C++, basicstyle=\ttfamily\small]
PetscSection section;
DMGetLocalSection(dm, &section);
VecSet(imposed, 0.0);
PetscScalar* arr;
VecGetArray(imposed, &arr);
for (const auto& [nid, u] : bc_list) {
    PetscInt pt = domain.node(nid).sieve_id.value();
    PetscInt offset;
    PetscSectionGetOffset(section, pt, &offset);
    arr[offset] = u[0]; arr[offset+1] = u[1]; arr[offset+2] = u[2];
}
VecRestoreArray(imposed, &arr);
\end{lstlisting}

This is both efficient and correct: the PetscSection offset for a
constrained vertex is the same position used by
\texttt{DMGlobalToLocal}/\texttt{VecAXPY} in the assembly callbacks.

\subsection{Material State Commit and Revert}

After convergence, the total displacement (free + imposed) is computed
and passed to each element's \texttt{commit\_material\_state()}:

\begin{lstlisting}[language=C++, basicstyle=\ttfamily\small]
DMGlobalToLocal(dm, U, INSERT_VALUES, u_local);
VecAXPY(u_local, 1.0, imposed_solution());
for (auto& elem : model->elements())
    elem.commit_material_state(u_local);
VecCopy(u_local, model->state_vector());
\end{lstlisting}

The state vector is also updated for VTK post-processing.
On divergence, \texttt{revert\_material\_state()} restores each
material point to its previously committed state.


% =========================================================================
\section{Summary}
\label{sec:nlevo-summary}

The \texttt{NonlinearSubModelEvolver} completes the multi-scale
pipeline by enabling persistent, state-accumulating nonlinear
continuum sub-models.  Combined with the \texttt{MultiscaleCoordinator}
and the higher-order prismatic domain builder, it allows a complete
workflow where a global beam--column frame analysis drives local
three-dimensional concrete sub-models through an entire earthquake,
capturing progressive cracking, damage accumulation, and stress
redistribution.

The L-shaped building example demonstrates the full pipeline:
3-component earthquake input, linear-to-nonlinear transition,
critical-element sub-model creation, step-by-step evolution with
persistent material state, and comprehensive VTK + CSV output for
visualisation and quantitative post-processing.
